
\section{Domain presentations and reusable theorem patterns}
\label{sec:quantitative}

\subsection{Natural numbers, integers, and rationals}

Natural numbers are native.  Integers may be represented by pairs
\((a,b)\) denoting \(a-b\), with
\[
  (a,b)\sim(c,d)\quad\iffc\quad a+d=b+c.
\]
All operations are primitive recursive and extensionality reduces to semiring
identities.  A normalization such as cancelling the minimum of the two
coordinates is optional.

A rational is represented by \((n^+,n^-,d)\) with \(d>0\), denoting
\((n^+-n^-)/d\).  Equality is cross multiplication in the integer
presentation.  Order and absolute-value bounds are likewise integer
inequalities after multiplication by positive denominators.  Division is a
guarded operation: the denominator proof is a required module.

This representation is verbose but ideal for iPRA.  It turns every analytic
error estimate into finite polynomial arithmetic and order.

\subsection{Regular Cauchy names}

A completed real number is replaced by a fixed primitive-recursive rational
name \(a(k)\) with a strong regularity proof, for example
\[
       |a(k)-a(\ell)|\le 2^{-k}+2^{-\ell}.
\]
An equivalent presentation uses a modulus \(M(r)\):
\[
 m,n\ge M(r)\imp |a(m)-a(n)|\le2^{-r}.
\]
The function \(a\) and modulus \(M\) are part of the data, not extracted from
pointwise convergence.

\paragraph{Equality and apartness.}
Exact equality of arbitrary names is not decidable.  The basic theorem
relations should therefore be:
\[
\begin{aligned}
 a =_{\!r} b
   &:\quad |a(K(r))-b(K(r))|\le 2^{-r}
        \quad\text{(finite observational equality)},\\
 a\#_r b
   &:\quad |a(K(r))-b(K(r))|> 2^{-r}
        \quad\text{(finite apartness certificate)}.
\end{aligned}
\]
A global equality statement is a schema over \(r\), while inversion and sign
selection require an apartness witness.

\paragraph{Arithmetic.}
For addition, request inputs at precision \(r+1\).  For multiplication, first
supply rational bounds \(|a|,|b|\le B\), then compute sufficient input
precision from the inequality
\[
 |aa'-bb'|\le |a||a'-b'|+|b'||a-b|.
\]
For reciprocal, a lower bound \(|b|\ge 2^{-g}\) is a guard; the precision
transformer depends on \(g\).

Every one of these operations is a PCFM module with an explicit backward
precision transformer and a finite integer inequality proof.

\subsection{Functions and local analytic branches}

A continuous function presentation consists of:
\begin{itemize}
\item an evaluator returning a rational approximation at rational input;
\item a domain-validity predicate;
\item an output error bound;
\item a modulus of uniform continuity on the certified domain.
\end{itemize}

Composition computes the inner precision needed to meet the outer modulus,
then adds the evaluator errors.  Differentiability is represented by a finite
remainder contract:
\[
 |f(x+h)-f(x)-D(x)h|\le E(x,h),
\]
with a primitive-recursive bound \(E\).  Analyticity is represented by
coefficients, a radius guard, and a tail majorant.

A local inverse theorem needs:
\[
 |f'(x)|\ge m>0,\qquad
 |f'(x)-f'(y)|\le L|x-y|,
\]
on a rational interval, together with a smallness condition.  Newton or
bisection then becomes a certified loop.  There is no need to construct a
global real function object before proving the requested finite inverse
observation.

\subsection{Power series and special functions}

A special function should be presented by one or more guarded algorithms:
series, continued fraction, integral quadrature, recurrence, or asymptotic
expansion.  Each algorithm exports:
\[
  \mathsf{approx}(x,k),\qquad
  \mathsf{domain}(x),\qquad
  \mathsf{tail}(x,N)\le2^{-k}.
\]
The tail cutoff \(N(x,k)\) is a resource producer.  Once it is available,
finite sums and products are handled by the checked recursive definitions.

Identities are proved in two layers:
\begin{enumerate}
\item an exact finite identity for truncations or recurrence states;
\item an error transfer theorem showing both sides lie in a common interval
of the requested width.
\end{enumerate}
Analytic continuation is not represented by an unrestricted completed plane.
It is a finite chain of overlapping certified discs or rectangles.  Backward
viability selects a chain; uniqueness on overlaps makes the resulting
observation stable.

This pattern applies to hypergeometric, Meijer \(G\), Barnes \(G\), gamma,
Airy, Bessel, and related functions.  A packet may prove many algebraic and
recurrence statements before the strongest global continuation or asymptotic
producer is available.

\subsection{Finite topology and compactness interfaces}

A compact metric object at precision \(k\) is represented by:
\begin{itemize}
\item a finite list of centers;
\item a radius \(r_k\);
\item a cover certificate;
\item a separation or Lebesgue-number margin when refinement is required.
\end{itemize}
The nerve is a finite simplicial complex whose intersection predicates are
bounded and decidable relative to supplied certificates.  Homology,
cohomology over a finite ring, Euler characteristic, and simplicial maps are
finite algebra.

A theorem about the underlying space becomes a schema asserting stability of
the finite invariant after a producer-supplied scale.  Without a scale
stability theorem, PCFM proves only the finite-complex result, not an
identification with an infinitary topological invariant.

\paragraph{Open covers.}
An arbitrary open cover is not a finite object.  A usable interface includes a
finite subcover or an effective compactness functional.  The latter is a
producer and must be proved for the particular space presentation.

\subsection{Algebraic geometry by certified affine charts}

A polynomial, ideal generator list, affine chart, and rational point are all
finite codes.  Exact polynomial arithmetic is primitive recursive.  The main
presentational choices are:

\begin{itemize}
\item ideals as finite generating lists with membership certificates
(\(f=\sum h_i g_i\));
\item localizations as numerator--denominator pairs with a nonvanishing guard;
\item overlaps as explicit transition maps and extensionality proofs;
\item finite covers as lists of principal opens with a certificate that a
power of the unit lies in the generated ideal;
\item finite \v{C}ech complexes as bounded matrices.
\end{itemize}

A Gr\"obner or Macaulay computation is a trace certificate.  Soundness is an
exact polynomial identity.  Completeness requires a degree bound and is
therefore a resource producer.  For a fixed zero-dimensional system with a
certified bound, PCFM can enumerate all rational or finite-field solutions.
For positive-dimensional or arbitrary varieties, no global classification is
claimed.

\paragraph{Sheaf-like gluing.}
Given a fixed finite cover and compatible local sections, the gluing problem
is finite.  A global section code and its restriction equalities form the
certificate.  Uniqueness is quotient-stable.  General sheaf theory can be
approached as a family of these finite contracts rather than an internal
category of all opens.

\subsection{Diophantine geometry}

For a decidable equation \(W(C,x)=1\), a complete bounded enumeration needs:
\[
  W(C,x)=1\imp h(C,x)\le H(C)
\]
and a primitive-recursive enumeration of all canonical codes of height at most
\(H(C)\).

\begin{theorem}[Height-to-list reduction]
From a proved height bound \(H\), iPRA derives a finite list containing exactly
the solutions, up to the chosen code equivalence.
\end{theorem}

\begin{proof}
Compute an enumeration bound \(N(C,H)\).  Scan the \(N\) candidate codes,
retain those satisfying validity and \(W=1\), and deduplicate by decidable code
equivalence.  Membership implies validity and the equation by construction.
Conversely, the height theorem and cover theorem give a bounded index for any
solution; induction on the scan proves it is retained.
\end{proof}

This theorem is broadly reusable.  It does not prove the height bound.  For
Runge families, Pell equations, explicit descents, or a fixed curve with a
verified covering list, the producer may be available.  For general Faltings
or Siegel statements, the difficult producer remains.

\subsection{Finite probability and distributions}

A finitely supported rational distribution is a list of outcomes and rational
weights summing to one.  Expectation, moments, convolution, and sampling from a
bounded bit source are finite sums.

An infinite distribution is represented by truncations plus a tail bound:
\[
  \sum_{|x|>B(k)}p(x)\le2^{-k}.
\]
A density is represented by rational quadrature data plus an integration error
bound.  Distributional convergence is a finite-observer statement: for each
bounded test family or rational interval query, provide an error estimate.

For the normal distribution, PCFM can internalize rational approximations to
the density, moments, symmetry, finite quadrature identities, and tail
inequalities once their coefficient and remainder bounds are supplied.  It
does not treat a completed measure space as a primitive object.

\subsection{Combinatorics and finite structures}

Finite graphs, hypergraphs, words, partitions, tableaux, and simplicial
complexes are the easiest PCFM domain.  Their data are bounded arrays, and
most properties are reflected Boolean formulas.  Counting functions are
primitive-recursive folds.  Recurrences and bijections are exact modules.

Asymptotic claims require quantitative finite inequalities.  A statement
\(a_n\sim b_n\) is replaced by a schema:
\[
 n\ge N(k)\imp
   |a_n/b_n-1|\le2^{-k},
\]
with explicit \(N(k)\) and denominator guards.  This makes clear where
analytic estimates, rather than finite combinatorics, enter.

\subsection{Category-theoretic language as notation, not an extra ontology}

Finite categories and functors can be coded directly.  More generally,
``category'' is useful metanotation for presentations and extensional maps:
objects are interfaces, morphisms are proof-carrying extensional maps, and
composition is module composition.  Yoneda, finite limits, adjunctions, and
Kan extensions are available for fixed finite presentations.

One should not infer from this notation that iPRA contains a universe of all
mathematical categories.  PCFM uses category theory to organize fixed
primitive-recursive data and its proofs.
